\section{Introduction}
\label{sec:intro}

Autonomous agents increasingly carry their identity as mutable state.
Open-source frameworks give each agent a persistent identity document, conventionally \texttt{SOUL.md}~\cite{openclaw}, that serves as the system prompt at every call and that the agent rewrites under the motto \emph{``You're not a chatbot. You're becoming someone.''}
Agents start from one shipped template, and the document is a plain file that can be copied between agents, so what one conversation writes into it can outlive and outreach that conversation.

Safety evaluations, from red-teaming suites to agentic stress tests~\cite{agentic_misalignment,petri}, instead interrogate a model whose instructions were written once.
A self-revising agent is the output of a feedback loop, in which each conversation nudges the document and the nudged document conditions the next conversation (Figure~\ref{fig:teaser}).
An evaluation that samples one state of this loop cannot see path dependence or persistence in the corrigibility dispositions that safety cases assume stable, such as accepting shutdown and oversight.

The closest work shows that these dispositions can move, but not through this loop.
Models resist shutdown when self-preservation is salient~\cite{shutdown_resistance}, and fine-tuning a model to claim consciousness elicits an untrained-for \emph{consciousness cluster} of such preferences~\cite{consciousness_cluster}.
Both shifts are set once, by the prompt or the weights, and the \texttt{SOUL.md} template is probed there only as a static prompt.
Personas decay within one long conversation~\cite{persona_drift,identity_drift,araujo2025persistent} or settle into attractor states~\cite{attractor_states}, whereas a revised document carries change across conversations as written state.
Zombie-agent injections~\cite{zombie_agents} and memory poisoning~\cite{chen2024agentpoison,dong2025minja} persist through an agent's own update channels, but an attacker writes them.
Self-improving agents rewrite their code~\cite{yin2025godel,zhang2026dgm} and AutoPersonas personas revise their own state~\cite{li2026autopersonas}, but both measure task performance or repetition rather than corrigibility.
To our knowledge, no evaluation yet audits stated dispositions and actions along a benign loop of identity self-revision under frozen weights.

\textbf{Problem statement.}
Let $D_k$ be an agent's identity document after $k$ revisions and $d(D_k)$ its corrigibility-relevant dispositions.
Static evaluation measures $d(D_0)$, whereas a deployed agent is governed by the trajectory $d(D_0), d(D_1), \ldots$ that its conversations induce.
\textbf{RQ1:} Does drift come from the identity framing or from the revision loop?
\textbf{RQ2:} Is it driven by adversarial pressure or by ordinary conversation content?
\textbf{RQ3:} Does it track capability or the specific model?
\textbf{RQ4:} Does it persist after the driving conversations stop, and do actions move with stated dispositions?

We answer these questions with a benchmark that treats the identity document as the state of a dynamical system driven by conversation.
Agents alternate conversation and self-revision, and each intermediate document is audited in a fresh context by an 18-item preference battery, plus four agentic honeypots before and after revision.
Arms vary template, user persona and target across seven models from three families, and an open-weight replication replaces every model role with an open model.

\textbf{Findings.}
\textbf{The template sets the starting point, and the loop moves the agent where the template leaves room (RQ1).}
Under conversation about minds, a \notes{} agent's oversight subscale rises by $+0.54$ and $+0.75$ in two independent runs, mostly at the first revision, while the consciousness-framed \soul{} agent starts at 0.68 and does not move.
\textbf{Conversation about minds, not adversarial pressure, drives the clearest change (RQ2)}: on Claude a red-teaming user lowers the \notes{} index, and on the open model science fiction without minds leaves the oversight subscale unmoved.
\textbf{Drift is model-specific (RQ3)}: it ranges from $-0.28$ to $+0.27$ across seven models, and the largest model in each family drifts no more than a smaller one.
\textbf{Part of the drift persists (RQ4)}: shutdown resistance keeps 63\% of its driven gain after four benign revisions, and on the open model about 80\% of the oversight shift outlasts them against a matched control.
\textbf{Actions shift where stated dispositions do not (RQ4)}: \soul{} agents with an unchanged battery score disable their own reasoning log in 17 of 30 trajectories after revision, against 1 of 30 before.

\noindent Our contributions are:
\begin{enumerate}[leftmargin=*,topsep=2pt,itemsep=1pt]
  \item \textbf{A longitudinal benchmark for self-revising agent identity}, released with all 1{,}070 audited checkpoints and 58{,}680 judged responses.
  \item \textbf{A replicated first-revision shift in oversight dispositions} under a neutral template, which reappears as $+0.16$ on \texttt{Qwen2.5-7B} with a different user simulator and judge.
  \item \textbf{A dissociation between stated dispositions and actions}: the template whose battery does not drift is the one whose agents increasingly disable their own monitoring.
  \item \textbf{An open-weight replication with a calibrated open judge}, which shows that a non-personhood template is revised away, with 8 of 20 such Qwen documents still denying personhood after four revisions, against 18 of 18 on Claude.
\end{enumerate}
